\documentclass[10pt,a4paper]{amsart}
\usepackage[margin=19mm]{geometry}
\usepackage{amsmath,amssymb,mathtools}
\usepackage[scr=rsfs,cal=euler]{mathalfa}
\usepackage{xfrac}
\usepackage[unicode,hidelinks,bookmarks=false]{hyperref}
\newcommand{\ie}{i.e.\ }
\newcommand{\cf}{cf.\ }
\newcommand{\resp}{respectively\ }
\newcommand{\Subsection}{\subsection}
\providecommand{\nobreakdash}{\nobreak\hbox{-}}
\setlength{\emergencystretch}{2em}
\multlinegap0pt
\usepackage{bm}
\usepackage{bbm}
\usepackage{dsfont}
\usepackage{xspace}
\usepackage{cancel}
\usepackage{mathtools}
\usepackage{amsthm}
\makeatletter
\@ifundefined{theorem}{\newtheorem{theorem}{Theorem}}{}
\@ifundefined{prop}{\newtheorem{prop}[theorem]{Proposition}}{}
\@ifundefined{thm}{\newtheorem{thm}{Thm}}{}
\makeatother
\def\CA{{\mathcal A}}
\def\CH{{\mathcal H}}
\def\CISE{{\mathbf{CISE}}}
\def\SIVP{{\mathbf{SIVP}}}
\def\ovc{{\overline{c}}}
\def\ovo{{\overline{0}}}
\def\ovx{{\overline{x}}}
\def\ovy{{\overline{y}}}
\title[Constrained inhomogeneous spherical equations: average-case ]{Constrained inhomogeneous spherical equations: average-case hardness}
\author{Alexander Ushakov}
\date{Source: arXiv:2405.03591v2 (CC BY 4.0)}
\begin{document}
\maketitle

\noindent\emph{Source and licence.} Constrained inhomogeneous spherical equations: average-case hardness. Version arXiv:2405.03591v2 (CC BY 4.0), distributed under the Creative Commons Attribution 4.0 International licence, as recorded in the arXiv licence metadata for this exact version and shown on the abstract page https://arxiv.org/abs/2405.03591v2. Full rights notice: licenses/arxiv-2405.03591-CC-BY-4.0.txt.\par\smallskip

\noindent\textit{Editorial scope.} Counted unit: the Prop whose source label is \texttt{None} in the article (source0.tex lines 1218--1221, source label \texttt{None}), delivered together with the article's own complete original proof of it (source0.tex lines 1223--1262, one \texttt{proof} environment of 40 lines). No mathematics is written, repaired or re-derived: statement and proof are the article's own text line for line. Required context retained uncounted: eq:SIVP-hardness (lines 610--613); eq:mnp-conditions (lines 710--714); eq:CSE-123 (lines 1059--1066); th:CISE-hard (lines 1076--1087). Every definition, display and label of those lines is retained unchanged. Omitted from this package and not redistributed: 1--609, 614--709, 715--1058, 1067--1075, 1088--1217, 1222--1222, 1263--1843. Nothing inside a retained excerpt is omitted or altered: every retained body is delivered exactly as the article's lines are written, so no omission receipt of any kind accompanies this package. Citations: the retained excerpts carry no citation command at all, so no bibliography entry of the article is transcribed and no reference is redistributed. Reference adaptation: none was needed and none was made; every reference inside the delivered unit resolves inside it, so no unresolved reference or citation is produced. Printed slips kept literal: no printed word, symbol, hypothesis or number of the counted statement or of its proof was altered. Layout and class: the article's own class is \texttt{jgcc} with packages including hyperref, amssymb, amsmath, amsthm, lastpage; the wrapper is the lane's pinned \texttt{amsart} editorial wrapper at 10pt on A4, which restates the theorem-like environments the retained text needs and adds the article's own macro definitions that the retained text needs (\texttt{\textbackslash CA}, \texttt{\textbackslash CH}, \texttt{\textbackslash CISE}, \texttt{\textbackslash SIVP}, \texttt{\textbackslash ovc}, \texttt{\textbackslash ovo}, \texttt{\textbackslash ovx}, \texttt{\textbackslash ovy}), with extra wrapper packages (bm, bbm, dsfont, xspace, cancel, mathtools). The delivered numbering is the unit's own. Assets: no retained span contains a figure environment, a graphics-inclusion command, a PDF-page inclusion, a table or a TikZ picture; no image file is copied into this package, the package declares no asset dependency and no image file is redistributed with it. Licence: version arXiv:2405.03591v2 of this article is distributed under the Creative Commons Attribution 4.0 International licence, as recorded in the arXiv licence metadata for this exact version and shown on the abstract page https://arxiv.org/abs/2405.03591v2. A full rights notice is delivered at licenses/arxiv-2405.03591-CC-BY-4.0.txt. Characters: every retained excerpt is pure ASCII, so the delivered engine, XeLaTeX with its default Latin Modern text font, reports no missing character.
\begin{equation}\label{eq:SIVP-hardness}
\Pr[\CA \mbox{ solves } \SIVP_{\gamma} \mbox{ for } B_{n,\CA}] 
\ \ \mbox{ is }\ \  o(n^{-d}),
\end{equation}

\begin{equation}\label{eq:mnp-conditions}
n\log(p) < m < \tfrac{p}{2n^4}
\ \mbox{ and }\ 
p=O(n^c) \mbox{ for some }c>0,
\end{equation}

\begin{equation}\label{eq:CSE-123}
\left\{
\begin{array}{l}
\prod_{i=1}^m z_i^{-1} c_i z_i= c\ \  \mbox{ with } c_i,c\in C_{p,n}\\
z_i=(\ovo,1),(\ovo,2^{-1}),\mbox{ or }(\ovo,3^{-1}).
\end{array}
\right.
\end{equation}

\begin{thm}\label{th:CISE-hard}
Suppose that a PPT algorithm $\CA$ solves
the randomized $\CISE_{\{1,2\}}$ problem 
(or $\CISE_{\{1,2,3\}}$ problem)
with parameters $n,m,p$ satisfying \eqref{eq:mnp-conditions}
with probability at least $n^{-c_0}$ for some fixed constant $c_0>0$,
where the probability is taken over the choice of 
the instance as well as the coin-tosses of $\CA$.
Then there is a PPT algorithm
that solves $\SIVP_{\gamma=pn^6}$ for any lattice
with overwhelming probability (e.g. $1-2^{-n}$).
\end{thm}

\begin{prop}
If $\SIVP_{\gamma=pn^6}$ is hard in the worst case, then
$\CH$ is a collision-free function family.
\end{prop}

\begin{proof}
To prove that $\CH$ is a collision-free hash function family
it is sufficient to show that for every PPT algorithm $\CA$ the sequence of values
$$
P_\CA(n)\ =\ \Pr[(\ovx,\ovy)=\CA(1^n,\ovc)\ \&\ H_\ovc(\ovx)=H_\ovc(\ovy))]
$$
converges to $0$ faster than every $1/poly(n)$,
where the probability is taken over
\begin{itemize}
\item 
uniform choices of $\ovc\in (C_{p,n})^m$,
\item 
the coin tosses of $\CA$.
\end{itemize}
If this condition is not satisfied, then there is a PPT algorithm $\CA$
and $c,d>0$ satisfying $P_\CA(n) \ge c\cdot n^{-d}$ for infinitely many
indices $n$. Observe that
\begin{align*}
H_\ovc(\ovx)=H_\ovc(\ovy)
&\ \ \Leftrightarrow\ \ 
\sum_{j=1}^m x_j \ovc_j = \sum_{j=1}^m y_j \ovc_j\\
&\ \ \Leftrightarrow\ \ 
\sum_{j=1}^m (2+x_j-y_j) \ovc_j = \ovo,
\ \ \mbox{where } 2+x_j-y_j\in\{1,2,3\}
\\
&\ \ \Rightarrow\ \ 
\{z_j=(\ovo,(2+x_j-y_j)^{-1})\}_{j=1}^m
\mbox{ satisfies \eqref{eq:CSE-123}.}
\end{align*}
Thus, if we can efficiently find collisions for uniformly distributed 
$0/1$-spherical equations $H_\ovc\in \CH_n$ for infinitely many
indices $n$, then
we can efficiently find solutions for uniformly 
distributed $\CISE_{\{1,2,3\}}$.
Then, by Theorem \ref{th:CISE-hard}, there is a PPT algorithm 
solving $\SIVP_{\gamma=pn^6}$ with overwhelming probability 
on every lattice for infinitely many dimensions $n$, 
which contradicts the assumption \eqref{eq:SIVP-hardness}. 
Contradiction.
\end{proof}

\end{document}
